% AI Email Triage Agent — Workflow Automation Case Study
% Compile: pdflatex case-study.tex

\documentclass[11pt,a4paper]{article}

\usepackage[margin=1in]{geometry}
\usepackage{graphicx}
\usepackage{hyperref}
\usepackage{enumitem}
\usepackage{booktabs}
\usepackage{xcolor}
\usepackage{parskip}

% Renders image if file exists; otherwise shows a labeled placeholder box.
\newcommand{\optionalimage}[2][0.92\textwidth]{%
  \IfFileExists{#2}{%
    \includegraphics[width=#1]{#2}%
  }{%
    \fbox{\parbox[c][4.5cm][c]{#1}{\centering\textit{Add \texttt{#2} to this folder}}}%
  }%
}

\newcommand{\screenshot}[3][0.9\textwidth]{%
  \begin{figure}[htbp]
    \centering
    \optionalimage[#1]{#2}
    \caption{#3}
  \end{figure}
}

\hypersetup{
  colorlinks=true,
  linkcolor=black,
  urlcolor=blue,
  citecolor=black
}

\title{\textbf{Workflow Automation Case Study}\\[0.4em]
\large AI Email Triage Agent for University Student Services}
\author{Asif Khan}
\date{June 2026}

\begin{document}

\maketitle

% ---------------------------------------------------------------------------
\section*{Try the Agent}
% ---------------------------------------------------------------------------

\begin{itemize}[leftmargin=*, nosep]
  \item \textbf{Live agent (email):} Send a message to \href{mailto:asifbot3715@gmail.com}{asifbot3715@gmail.com}.
        The agent will classify your request and reply automatically, ask for missing information, or escalate to a human when needed.
        Allow up to one minute for a response (the inbox is polled every minute).
  \item \textbf{Operations dashboard:} \url{<PASTE\_YOUR\_DASHBOARD\_URL\_HERE>}
\end{itemize}

\vspace{0.5em}
\noindent\textit{Note: This project was built in a single day before the 3 June 2026 deadline. Scope was intentionally kept simple: a working end-to-end agent with deployment, audit logging, and escalation---not a production-hardened system.}

% ---------------------------------------------------------------------------
\section{Why This Problem?}
% ---------------------------------------------------------------------------

University student-services teams receive a high volume of repetitive email: transcript requests, fee questions, registration issues, password resets, and general inquiries. Much of this work is \emph{data-heavy} (student ID, course codes, deadlines) yet still handled manually---staff read each message, decide urgency, reply or forward, and log outcomes in spreadsheets or ticketing tools. That repetition consumes real time and delays students who need quick answers.

This maps directly to the \textbf{Intake Triage Bot} pattern from the brief: classify incoming requests, gather missing information through conversation, and route complex cases to the right person with full context.

% ---------------------------------------------------------------------------
\section{Why I Picked This Workflow}
% ---------------------------------------------------------------------------

With the application deadline on 3 June 2026, I had roughly \textbf{one day} to design, build, test, and deploy a working agent. I chose email triage because:

\begin{itemize}[nosep]
  \item It is a \textbf{real} workflow I have seen in education settings, not a toy demo.
  \item The user interaction model is familiar (email in, email out) and easy for reviewers to test.
  \item I could ship an autonomous loop without renting a server, by using \textbf{Google Apps Script} for Gmail, scheduling, and API access.
  \item The branching logic (reply / clarify / escalate) is enough to show agentic behavior within a tight time budget.
\end{itemize}

% ---------------------------------------------------------------------------
\section{Who Is the User?}
% ---------------------------------------------------------------------------

\textbf{Primary user:} A student emailing the student-services inbox with a question or request.

\textbf{Secondary user:} A staff member or intern monitoring the queue on a \textbf{live dashboard} and handling \textbf{escalated} cases forwarded by email.

\textbf{When and how:} Students send email at any time; the agent runs on a one-minute schedule. Staff use the dashboard during office hours to see what was processed, what failed, and what needs human follow-up.

% ---------------------------------------------------------------------------
\section{Architecture}
% ---------------------------------------------------------------------------

\subsection*{How the Agent Works}

\begin{enumerate}[nosep]
  \item A \textbf{time trigger} (every 1 minute) runs \texttt{pollInbox} in Google Apps Script.
  \item Unread inbox messages are fetched (with deduplication by Gmail message ID and thread labels).
  \item Message content and \textbf{thread history} are sent to \textbf{Groq} (\texttt{llama3-70b-8192}) with a fixed JSON schema.
  \item The agent chooses an action and executes it via Gmail; every step is logged to \textbf{Google Sheets}.
  \item A \textbf{React dashboard} polls the Apps Script web app every 5 seconds for live status.
\end{enumerate}

\subsection*{Autonomous Decisions vs.\ Escalation}

\begin{center}
\begin{tabular}{@{}ll@{}}
\toprule
\textbf{Decision} & \textbf{When} \\
\midrule
\texttt{auto\_reply} & High-confidence, routine inquiry; agent sends a direct answer \\
\texttt{clarify}     & Required information missing (e.g.\ student ID); agent asks one question; thread stays open \\
\texttt{escalate}    & Low confidence, urgent keywords, or out-of-scope; student gets acknowledgment; staff notified \\
\bottomrule
\end{tabular}
\endcenter

\subsection*{Failure Handling}

\begin{itemize}[nosep]
  \item Invalid or failed Groq responses $\rightarrow$ email marked \texttt{failed}, generic fallback reply sent, thread labeled to avoid infinite retries.
  \item Gmail daily send quota exceeded $\rightarrow$ processing stops gracefully; errors visible in Apps Script execution logs.
  \item Ambiguous or off-topic input $\rightarrow$ model prompted to escalate rather than guess.
  \item Multi-turn clarify flows $\rightarrow$ follow-up replies in the same thread are processed with full conversation context.
\end{itemize}

\subsection*{Architecture Diagram}

\begin{figure}[htbp]
  \centering
  \optionalimage{arch.png}
  \caption{System architecture: Gmail $\rightarrow$ Apps Script agent $\rightarrow$ Groq / Sheets / replies; React dashboard reads from Apps Script API.}
  \label{fig:architecture}
\end{figure}

\subsection*{Technology Stack}

\begin{center}
\begin{tabular}{@{}ll@{}}
\toprule
\textbf{Layer} & \textbf{Technology} \\
\midrule
Ingestion \& actions & Google Apps Script (\texttt{GmailApp}) \\
LLM classification  & Groq API (\texttt{llama3-70b-8192}) \\
Audit log           & Google Sheets \\
Monitoring UI       & Vite, React, Tailwind CSS \\
\bottomrule
\end{tabular}
\end{center}

% ---------------------------------------------------------------------------
\section{Screenshots}
% ---------------------------------------------------------------------------

\noindent Place \texttt{sc1.png}--\texttt{sc4.png} in the same folder as this \texttt{.tex} file before compiling.

\screenshot{sc1.png}{Dashboard overview: live queue, status counts, and email list.}

\screenshot{sc2.png}{Email detail modal: intent, confidence, reply sent, and extracted entities.}

\screenshot{sc3.png}{Gmail thread showing agent auto-reply, clarify question, or escalation acknowledgment.}

\screenshot{sc4.png}{Google Sheets audit log with processed rows and Groq response data.}

% ---------------------------------------------------------------------------
\section{What I Learned}
% ---------------------------------------------------------------------------

\begin{itemize}[nosep]
  \item \textbf{Agent design is state + policy:} Labeling threads, tracking message IDs, and keeping clarify threads open matter as much as the prompt.
  \item \textbf{Serverless can ship fast:} Apps Script removed the need for a backend deployment while still supporting Gmail actions and a JSON API.
  \item \textbf{Quotas are real constraints:} Gmail send limits and Groq rate limits forced better deduplication and careful testing.
  \item \textbf{Escalation is a feature:} Knowing when \emph{not} to auto-reply is what makes the agent trustworthy for a demo inbox.
  \item \textbf{Observability helps demos:} A live dashboard plus a Sheet audit trail made debugging and presentation much easier under time pressure.
\end{itemize}

\end{document}
